\documentclass[a4paper, 10pt, ]{article}

\usepackage[slovak]{babel}

% ------------------------------

\usepackage[utf8]{inputenc}
\usepackage[T1]{fontenc}

\usepackage[left=4cm,
            right=4cm,
            top=2.1cm,
            bottom=2.6cm,
            footskip=7.5mm,
            twoside,
            marginparwidth=3.0cm,
            %showframe,
            ]{geometry}

\usepackage{graphicx}
\usepackage[dvipsnames]{xcolor}
% https://en.wikibooks.org/wiki/LaTeX/Colors

% ------------------------------

\usepackage{lmodern}

\usepackage[tt={oldstyle=false,proportional=true,monowidth}]{cfr-lm}
% https://mirror.szerverem.hu/ctan/fonts/cfr-lm/doc/cfr-lm.pdf

% ------------------------------

\usepackage{amsmath}
\usepackage{amssymb}
\usepackage{amsthm}

\usepackage{booktabs}
\usepackage{multirow}
\usepackage{array}
\usepackage{dcolumn}

\usepackage{natbib}

% ------------------------------

\hyphenpenalty=6000
\tolerance=1000

\def\naT{\mathsf{T}}

% ------------------------------

\makeatletter

    \def\@seccntformat#1{\protect\makebox[0pt][r]{\csname the#1\endcsname\hspace{4mm}}}

    \def\cleardoublepage{\clearpage\if@twoside \ifodd\c@page\else
    \hbox{}
    \vspace*{\fill}
    \begin{center}
    \phantom{}
    \end{center}
    \vspace{\fill}
    \thispagestyle{empty}
    \newpage
    \if@twocolumn\hbox{}\newpage\fi\fi\fi}

    \newcommand\figcaption{\def\@captype{figure}\caption}
    \newcommand\tabcaption{\def\@captype{table}\caption}

\makeatother

% ------------------------------

\usepackage{fancyhdr}
\fancypagestyle{plain}{%
\fancyhf{} % clear all header and footer fields
% \fancyfoot[C]{\sffamily {\bfseries \thepage}\ | {\scriptsize\oznacenieCasti}}
\fancyfoot[C]{\sffamily {\bfseries \thepage}{\color{Gray}\scriptsize$\,$z$\,$\pageref{LastPage}}\ | \includegraphics[height=5pt]{../../COMMONFILES/KUT_logo_v0.1.pdf}{\scriptsize\KUTporadoveCislo}}
\renewcommand{\headrulewidth}{0pt}
\renewcommand{\footrulewidth}{0pt}}
\pagestyle{plain}

% ------------------------------

\usepackage{titlesec}
\titleformat{\paragraph}[hang]{\sffamily  \bfseries}{}{0pt}{}
\titlespacing*{\paragraph}{0mm}{3mm}{1mm}
\titlespacing*{\subparagraph}{0mm}{3mm}{1mm}

\titleformat*{\section}{\sffamily\Large\bfseries}
\titleformat*{\subsection}{\sffamily\large\bfseries}
\titleformat*{\subsubsection}{\sffamily\normalsize\bfseries}


% ------------------------------

\PassOptionsToPackage{hyphens}{url}
\usepackage[pdfauthor={},
            pdftitle={},
            pdfsubject={},
            pdfkeywords={},
            % hidelinks,
            colorlinks=false,
            breaklinks,
            ]{hyperref}


% ------------------------------

\graphicspath{%
{../fig_standalone/}%
{../../PY/fig/}%
{../../ML/fig/}%
{./fig/}%
}

% ------------------------------

\usepackage{enumitem}

\usepackage{lettrine}

% ------------------------------

\usepackage{lastpage}

\usepackage{microtype}

% ------------------------------

\usepackage{algorithm}
\usepackage[noend]{algpseudocode}
\makeatletter
\renewcommand{\ALG@name}{Algoritmus}
\makeatother
\usepackage{amsmath}
\usepackage{bbold}
\usepackage{calc}
\usepackage{dsfont}
\usepackage{mathtools}
\usepackage{tabto}


\newcommand{\mr}[1]{\mathrm{#1}}
\newcommand{\bs}[1]{\boldsymbol{#1}}
\newcommand{\bm}[1]{\mathbf{#1}}

\newcommand{\diff}[2]{\frac{\Delta #1}{\Delta #2}}
\newcommand{\der}[2]{\frac{d #1}{d #2}}
\newcommand{\parder}[2]{\frac{\partial #1}{\partial #2}}

\newcommand{\argmax}[0]{\mr{argmax}}
\newcommand{\diag}[0]{\mr{diag}}
\newcommand{\rank}[0]{\mr{rank}}
\newcommand{\trace}[0]{\mr{tr}}

\renewcommand{\Re}{\mr{Re}}
\renewcommand{\Im}{\mr{Im}}


\theoremstyle{definition}
\newtheorem{definition}{Definícia}[section]
\newtheorem{theorem}{Veta}[section]
\newtheorem{lemma}[theorem]{Lemma}
\newtheorem{example}{Príklad}[section]
\renewcommand*{\proofname}{Dôkaz}

% ------------------------------

\usepackage{siunitx}
\usepackage{xcolor}

% Define a custom color
\definecolor{backcolour}{rgb}{0.95,0.95,0.92}
\definecolor{codegreen}{rgb}{0,0.6,0}

% Define a custom style
\lstdefinestyle{myStyle}{
    belowcaptionskip=1\baselineskip,
    breaklines=true,
    frame=single,
    showspaces=false,                % show spaces everywhere adding particular underscores; it overrides 'showstringspaces'
    showstringspaces=false,          % underline spaces within strings only
    showtabs=false,
    numbers=left, 
    basicstyle=\footnotesize\ttfamily,
    keywordstyle=\bfseries\color{blue!70!black},
    commentstyle=\color{green!40!black},
    identifierstyle=\color{gray!50!black},
    backgroundcolor=\color{white},
}

% Use \lstset to make myStyle the global default
\lstset{style=myStyle}

% -----------------------------------------------------------------------------

\def\oznacenieCelku{Kolekcia učebných textov}

% -----------------------------------------------------------------------------


\def\KUTporadoveCislo{devAS2}

\def\oznacenieVerzie{v1.0}
% \def\oznacenieVerzie{\phantom{v1.0}}

\def\mesiacRok{september 2025}

\def\authorslabel{RJ}






% -----------------------------------------------------------------------------

\begin{document}

% -----------------------------------------------------------------------------
% Uvodny nadpis

\noindent
\parbox[t][18mm][c]{0.3\textwidth}{%
    \raisebox{-0.9\height}{%
        \phantom{.}\includegraphics[height=18mm]{./COMMONFILES/URKFEIlogo.pdf}%
    }%
}%
\parbox[t][18mm][c]{0.7\textwidth}{%
    \raggedleft

    \sffamily
    \fontsize{16pt}{18pt}
    \fontseries{sbc}
    \selectfont

    \noindent
    \textcolor[rgb]{0.75, 0.75, 0.75}{\textls[25]{\oznacenieCelku}}
}%

\noindent
\parbox[t][16mm][b]{0.5\textwidth}{%
    \raggedright

    \color{Gray}
    \sffamily

    \fontsize{12pt}{12pt}
    \selectfont
    \mesiacRok

    \fontsize{6pt}{10pt}
    \selectfont
    \href{https://github.com/OkoliePracovnehoBodu/KUT}{github.com/OkoliePracovnehoBodu/KUT}

    \fontsize{8pt}{10pt}
    \selectfont
    \authorslabel




}%
\parbox[t][16mm][b]{0.5\textwidth}{%
    \raggedleft

    \sffamily

    \fontsize{6pt}{6pt}
    \selectfont

    \textcolor[rgb]{0.68, 0.68, 0.68}{\oznacenieVerzie}


    \fontsize{14pt}{14pt}
    \selectfont

    \bfseries

    \includegraphics[height=12pt]{./COMMONFILES/KUT_logo_v0.1.pdf}%
    {%
        \textls[-50]{\KUTporadoveCislo}
    }%
}%

% -----------------------------------------------------------------------------




\vspace{6mm}

% ---------------------------------------------
\sffamily
\bfseries
\fontsize{18pt}{21pt}
\selectfont

\begin{flushleft}
    Laboratórne zariadenie AeroShield:\\ firmwáre PIO
\end{flushleft}

\bigskip

% -----------------------------------------------------------------------------
\normalsize
\normalfont
% -----------------------------------------------------------------------------

\lstset{style=mystyle}










\noindent
\lettrine[lines=1, nindent=1pt, loversize=0.0]{C}{ieľom}
textu je opis firmwarovej stránky \emph{AS} implementovaná na Arduino UNO R3, ktorá slúži ako komunikačný most medzi Matlab-om a fyzickým zariadením \emph{AS}.


\section{Pomocné funkcie časovačov}
%\lstinputlisting[caption=Sample Code Listing C++, label={lst:listing-cpp}, language=C++]{../PIO/src/main.cpp}
\begin{lstlisting}[caption=Obslužné funkcie interných časovačov Arduina., label={code:timer_misc}, language=C++]
// -------------------------------
// Define Timer related macros, functions
// -------------------------------

#define CPUF 16000000UL                                      // CPU frequency for Arduino UNO
#define DIF(Tms) (1000.0f / Tms)                             // Frequency divider macro
#define CMR(PS, DIF) (static_cast<int>(CPUF / PS / DIF) - 1) // Compare match register value calculation macro
#define PS(CMR, DIR) ((CPUF / (CMR + 1)) / DIR)              // Prescaler value calculation macro
#define CHECK_TIMER(CMR) (CMR < 256)                         // Check if CMR value fits in 8-bit register
#define CURRENT_PS 1024 // Current prescaler value

// -------------------------------
// -------------------------------
\end{lstlisting}

Definícia premenných a pomocných funkcií na prácu a nastavenie interných časovačov Arduina - viď blok \ref{code:timer_misc}.

\section{Premenné zabudovanej LED}

\begin{lstlisting}[caption=Definície premenných zabudovanej LED., label={code:led_config}, language=C++]
// -------------------------------
// Built-in LED configs
// -------------------------------

const static int BUILT_IN_LED_PIN = 13; // PIN pre zabudovanu LED
const static unsigned long T_sample = 1000;
const static unsigned long LED_onTime = 15;
bool LED_on = false;

// -------------------------------
// -------------------------------
\end{lstlisting}

Definície premenných na zabezpečenie blikania v kontexte zabudovanej LED na zadefinovanej frekvencii - viď výpis \ref{code:led_config}, priradenie konštánt pre zjednodušenie následného použitia premenných v kóde. \texttt{LED\_on} premenná zabezpečuje prepínanie stavu LED medzi emitovaním a zhasnutým stavom.

\section{Premenné signálov a komunikácie}

\begin{lstlisting}[caption=Definície premenných signálov a komunikácie., label={code:sig_config}, language=C++]
// -------------------------------
// Required variables and constexpr functions
// -------------------------------

byte Ts_user = 1; // sampling time in ms

constexpr float TS_INTERVAL_MIN(byte Ts)
{
    return (static_cast<float>(Ts) * 1000.0f) * 0.95f;
}

float uSig = 0.0f, ySig = 0.0f, potSig = 0.0f, dtoffset = 0.0f, dtoffset_interval = 0.0f; // control, output, reference, time offset

unsigned long currentTime, dt, lastTime;

// -------------------------------
// -------------------------------
\end{lstlisting}

Definícia podstatných premenných, ktoré používame na uloženie terajších meraní, konfigurácia časovaní posielania dát, nakoniec ukladanie času trvania merania pre Arduino a diferencia času medzi terajšou komunikáciou a predošlou, môžeme vidieť na výpise kódu \ref{code:sig_config}. Hodnotu diferencie času sa používa na časovú synchronizáciu komunikácie, aby sa perióda vzorkovania hýbala v okolí žiadanej periódy vzorkovania, nakoľko zabezpečenie presnej periódy vzorkovania si vyžaduje hardwarové riešenie.

\section{Metódy čítania a zapisovania}

\begin{lstlisting}[caption=Obslužné funkcie čítania a zapisovania signálov., label={code:sig_funcs}, language=C++]
// -------------------------------
// API calls
// -------------------------------

void updateOutput()
{
    ySig = AeroShield.sensorReadDegree();
}

void updatePotentiometer()
{
    potSig = AeroShield.referenceRead();
}

void updateControl(const float &u)
{
    uSig = u;
}

void writeControl()
{
    AeroShield.actuatorWrite(uSig);
}

// -------------------------------
// -------------------------------
\end{lstlisting}

Funkcie, ktoré možeme vidieť na výpise kódu \ref{code:sig_funcs} slúžia iba na prehladnosť a čistosť kódu. Zabezpečujú zapisovanie hodnôt do definovaných premenných signálov alebo fyzikálnu realizáciu žiadanej akčnej veličiny v tomto prípade funkcia \texttt{writeControl()}.

\section{Sériová komunikácia}

\begin{lstlisting}[caption=Funkcia na posielanie dát po sériovej linke., label={code:send}, language=C++]
// -------------------------------
// Write data into serial comms
// -------------------------------

void printData()
{
    dt = currentTime - lastTime;
    dtoffset = dtoffset_interval - dt;
    if (dtoffset > 0)
    {
        if (dtoffset > 1000)
        {
            dtoffset = static_cast<unsigned long>(floor(dtoffset / 1000)); // round to nearest 1ms
            delay(dtoffset);
        }
        else
        {
            dtoffset = static_cast<unsigned long>(floor(dtoffset / 100) * 100); // round to nearest 100us
            delayMicroseconds(dtoffset);
        }
        currentTime = micros();
        dt = currentTime - lastTime;
    }
    lastTime = currentTime;

    Serial.print(currentTime);
    Serial.print(" ");
    Serial.print(uSig);
    Serial.print(" ");
    Serial.print(ySig);
    Serial.print(" ");
    Serial.print(potSig);
    Serial.print(" ");
    Serial.print(dt);
    Serial.println();
}

// -------------------------------
// -------------------------------
\end{lstlisting}

Kód \ref{code:send} je komunikačná funkcia, ktorá zabezpečuje posielanie meraných veličín po sériovej linke von - do Matlab-u. Zároveň implementuje algoritmus, ktorý zabezpečuje komunikáciu s kvázi stabilnou periódou vzorkovania a odosiela dáta čo najskôr. Nakoľko dĺžka trvania jedného cyklu \texttt{loop()} je variabilná a závisí od tohto času následná komunikácia, tak sme časovanie implementovali spôsobom: ak program môže poslať dáta skôr ako je žiadané, tak čaká \texttt{dtoffset} čas a následne posiela dáta, v prípade, že je neskôr, tak okamžite posiela dáta, ktoré sú najnovšie v pamäti.

\section{Čítanie signálov, sériovej linky a logika riadenia merania}
\begin{lstlisting}[caption=Čítanie signálov, sériovej linky a definícia poradia akcií riadenia., label={code:read}, language=C++]
// -------------------------------
// Read new values from serial com and inputs
// -------------------------------
void readData()
{
    updateOutput();
    updatePotentiometer();
}

int recvByte()
{
    if (Serial.available() >= 4)
    {
        Serial.readBytes((char *)&uSig, 4); // Read 4 bytes from Serial
        return 0;                           // Success
    }
    return 1; // No new data received
}

void processNewData()
{
    if (recvByte())
    {
        currentTime = micros();

        readData(); // Read the sensor and reference values

        writeControl(); // Write the control signal to the actuator

        printData();
    }
}
// -------------------------------
// -------------------------------
\end{lstlisting}

V bloku kódu \ref{code:read} sú 3 rôzne funkcie, z čoho sú 2 na čítanie dát a posledna na riadenia komunikácie medzi Matlabom, a AeroShield-om. V prvej funkcii \texttt{readData()} čítame hodnoty zo samotného AeroShield-u, teda fyzikálne veličiny ako natočenie vrtuľky a potenciometra. Funkcia \texttt{recvByte()} zabezpečuje čítanie žiadanej akčnej veličiny z Matlab-u po sériovej linke, čítanie je priamo implementované pre dátový typ \texttt{float32}. Nakoniec funkcia \texttt{processNewData()} zabezpečuje, meranie veličín zariadenia, spracovanie požiadaviek z Matlab-u a následnú fyzikálnu realizáciu - PWM signál posielaný priamo na \emph{Gate} pin MOSFET tranzistoru a odoslanie nameraných dát späť po sériovej linke.

\section{Inicializácia Arduina}

\begin{lstlisting}[caption=Inicializačná metóda Arduina., label={code:setup}, language=C++]
// -------------------------------
// Setup the Arduino board
// -------------------------------
void setup()
{
    Serial.begin(115200); // Initialize Serial communication at 115200 baud rate
    Serial.flush();
    pinMode(BUILT_IN_LED_PIN, OUTPUT); // for LED

    AeroShield.begin();
    AeroShield.calibrate();

    Serial.println("--- MCU config ---");
    while (!Serial.available())
    {
        delay(1); // Wait for user to open terminal
    }

    Ts_user = Serial.read(); // sampling time in ms

    if (Ts_user == 0)
    {
        Ts_user = 1; // minimum 1 ms
    }

    if (Ts_user > 10)
    {
        dtoffset_interval = static_cast<float>(Ts_user) * 1000.0f;
    }
    else
    {
        dtoffset_interval = TS_INTERVAL_MIN(Ts_user);
    }

    // Calculate the timers
    const int cmr = CMR(CURRENT_PS, DIF(T_sample));

    cli();

    // Timer1 for LED blink
    TCCR1A = 0; // Normal mode
    TCCR1B = 0;
    TCNT1 = 0;                           // Initialize counter value to 0
    OCR1A = cmr;                         // Set compare match register for desired timer count
    TCCR1B |= (1 << WGM12);              // CTC mode
    TCCR1B |= (1 << CS12) | (1 << CS10); // Set prescaler to 1024 and start the timer
    TIMSK1 |= (1 << OCIE1A);             // Enable timer compare interrupt

    sei();

    Serial.println("--- MCU starting [" + String(Ts_user) + "] ---");
    currentTime = micros();
    lastTime = currentTime - (Ts_user * 1000); // better for statistics

    // Initial reads and writes
    readData();
    updateControl(0.0f);
    writeControl();

    printData();
}
// -------------------------------
// -------------------------------
\end{lstlisting}

V rámci inicializačnej funkcie \ref{code:setup} sa nastavuje komunikácia po sériovej linke (Universal Serial Bus - \emph{USB}), módy GPIO pinov, inicializácia AeroShield firmwáru, perióda vzorkovania - posiela Matlab v našom prípade, výpočet doveleného oskorenia/oneskorenia posielania dát - slúži na zabezpečenie najrýchlejšej možnej komunikácie pre žiadanu periódu vzorkovania, časovač preblikávania zabudovanej LED, po celkovej inicializácii posielame správu o začiatku merania a nakoniec posielame počiatočný stav zariadena. Následne po poslaní počiatočného stavu sa začne volať funkcia \texttt{loop()} samotným Arduinom, ktorá následne riadi dej merania.

\section{Obslužná funkcia časovaču (spínanie LED)}
\begin{lstlisting}[caption=Spínanie LED v obslužnej funkcii \emph{Časovača 1.}, label={code:timer_callback}, language=C++]
// -------------------------------
// Timer 1 callback function for when triggered
// -------------------------------
ISR(TIMER1_COMPA_vect)
{
    LED_on = !LED_on;
    digitalWrite(BUILT_IN_LED_PIN, LED_on);
}
// -------------------------------
// -------------------------------
\end{lstlisting}

Metóda zobrazená vo výpise \ref{code:timer_callback} je špecialne definovaná samotným Arduinom, na ktorú sme sa pripojili, aby sa následne volal kód, ktorý sme vo vnútri tejto funkcie definovali. Funkcia je volaná iba vtedy, keď sa Časovač 1 zopne, podľa toho ako si definujeme interval zopínania časovača v inicializačnej funkcii.

\section{Hlavná slučka firmwáru}
\begin{lstlisting}[caption=Hlavná slučka Arduino programu., label={code:loop}, language=C++]
void loop()
{
    processNewData();
}
\end{lstlisting}

Funkciu znázornenú vo výpise kódu \ref{code:loop} možno nazvať hlavnou funkciou celého programu, nakoľko pomocou tejto funkcie sa Arduino rozhoduje, čo má v každom čase po inicializácii robiť. Funkciu \texttt{loop()} si možeme predstaviť ako nekonečnú slučku, ktorú Arduino volá vždy po inicializácii. To znamená, že v tejto funkcii definujeme, čo má Arduino robiť. V našom prípade sa volá funkcia \texttt{processNewData()}, ktorá sa stará o riadenie komunikácie medzi sériovou linkou a výstup/vstupom Arduina.


% -----------------------------------------------------------------------------

\end{document}

% -----------------------------------------------------------------------------